% answers_short.tex —— 1989 年·上海卷 · 简短答案速查
% 由 tools/make_short_answers.py 自 answers.tex 的 \gksolhead/\gkanswer 提取，
% 并与 1989/1989_数学真题_上海_答案要点.md 逐题比对后生成，勿手改（改源 MD 后重跑）。

\begin{gknotice}
\textbf{使用说明}：本节只给最终答案便于快速核对，完整推导见解析版。答案为编写者独立推导，\textbf{非官方标准答案}。
\end{gknotice}

\gkpart{大题}
\begin{itemize}[leftmargin=2.6em]
  \item[\textbf{1.}] (1) $\dfrac13$；(2) $-\sqrt2 i$；(3) $x=4k\pi+\dfrac{2\pi}3$ 或 $x=4k\pi-2\pi$，$k\in\mathbb Z$；(4) $\dfrac{16}3\text{ cm}^3$；(5) $(-\infty,0)$；(6) $\dfrac{60}{13}$；(7) $x=4$；(8) $-\dfrac{2\sqrt3}3$；(9) $\dfrac32$；(10) $y=-\dfrac{17}{16}$
  \item[\textbf{2.}] (1) C；(2) A；(3) B；(4) D；(5) B；(6) C；(7) D；(8) A；(9) B；(10) D
  \item[\textbf{3.}] (1) $-\dfrac{5+6\sqrt2}{14}$；(2) $(0,\dfrac\pi2)\cup[\pi,4]$；(3) $x+7y-25=0$ 或 $7x-y+25=0$；(4) $r=4$；$T_{16}=-15504x^{30}$，$T_6=-15504x^{10}$
  \item[\textbf{4.}] 见解析版
  \item[\textbf{5.}] (1) 线段 $AB$，端点为 $A(2+\sqrt2,\sqrt2-1)$、$B(2-\sqrt2,-\sqrt2-1)$；(2) $[2\pi-\operatorname{arctg}\dfrac{3\sqrt2+4}{2},2\pi)\cup[0,\operatorname{arctg}\dfrac{3\sqrt2-4}{2}]$；(3) $[\dfrac{3\sqrt2}{2},\sqrt{9+2\sqrt2}]$
  \item[\textbf{6.}] 令 $\varphi=\dfrac{\sqrt5-1}{2}$；$-1<q<\varphi$ 且 $q\ne0$ 时 $A_n>B_n$，$q=\varphi$ 时 $A_n=B_n$，$q>\varphi$ 时 $A_n<B_n$。
  \item[\textbf{7.}] (1) $\rho=\dfrac1{1-\frac1p\cos\theta}$，其中 $0<\cos\theta<p$；$0<p<1$ 为双曲线部分，$p=1$ 为抛物线部分，$p>1$ 为椭圆部分；(2) $4$（$0<p\le2$）；$\dfrac{p^2}{p-1}$（$p>2$）
\end{itemize}
